\documentclass[a4paper,10pt]{article}

\usepackage[utf8]{inputenc}
\usepackage[T1]{fontenc}
\usepackage{geometry}
\usepackage{booktabs}
\usepackage{array}
\usepackage{tabularx}
\usepackage{longtable}
\usepackage{listings}
\usepackage{xcolor}
\usepackage[hidelinks]{hyperref}
\usepackage{amsmath}


\geometry{
    a4paper,
    top=1.8cm,
    bottom=1.8cm,
    left=1.8cm,
    right=1.8cm
}

\setlength{\parindent}{0pt}
\setlength{\parskip}{5pt}
\renewcommand{\arraystretch}{1.18}


\newcommand{\hex}[1]{\texttt{0x#1}}
\newcommand{\code}[1]{\texttt{#1}}

\title{\textbf{DEOS Brake Node \& Service ICD}\\
\large Sürüm 1.1 -- Fren Düğümü (\hex{12}) ve Servisi (\hex{03})}
\author{\textbf{Emir Arkalı} (\href{mailto:emirarkali07@gmail.com}{emirarkali07@gmail.com})\\DEOS -- Dokuz Eylül Otonom Sistemler}
\date{İlk Yayın: 28 Eylül 2026 \quad | \quad Son Revizyon: 13 Ekim 2026}

\begin{document}
\maketitle

\begin{center}
\small
\begin{tabularx}{\textwidth}{ll lX}
\toprule
\textbf{Doküman Tipi:} & Interface Control Document (ICD) & \textbf{Doküman Kodu:} & \code{DOC-ICD-NODE-BRK-010} \\
\textbf{Doküman Sürümü:} & Sürüm 1.1 & \textbf{Protokol Sürümü:} & v1.1 (\hex{11}) -- Medya Bağımsız \\
\textbf{Yazar:} & Emir Arkalı (\texttt{emirarkali07@gmail.com}) & \textbf{Son Değiştiren:} & Emir Arkalı \\
\textbf{İlk Yayın Tarihi:} & 28 Eylül 2026 & \textbf{Son Revizyon Tarihi:} & 13 Ekim 2026 \\
\textbf{Durum:} & Yürürlükte / Onaylandı & \textbf{Taşıma / Kapsam:} & Brake Düğümü (\hex{12}) \& BRAKE Servisi (\hex{03}) \\
\bottomrule
\end{tabularx}
\end{center}

\subsection*{Doküman Revizyon Geçmişi}
\begin{center}
\small
\begin{tabularx}{\textwidth}{c c l X}
\toprule
\textbf{Revizyon} & \textbf{Tarih} & \textbf{Yazar / Değiştiren} & \textbf{Açıklama} \\
\midrule
v0.3 & 28 Eylül 2026 & Emir Arkalı & İlk taslak ve monolitik ICD v0.3 tanımı. \\
v1.0 & 28 Eylül 2026 & Emir Arkalı & Fren alt sistemi bağımsız doküman revizyonu. \\
v1.1 & 13 Ekim 2026 & Emir Arkalı & Sürüm 1.1 standardizasyonu, düğüm ve servis ayrımı netleştirildi, kurumsal üstveri eklendi. \\
\bottomrule
\end{tabularx}
\end{center}
\vspace{0.8em}

\section{Fren (Brake) Düğüm ve Servis Mimarisi}

DEOS mimarisinde fren alt sistemi, hem donanımsal/mantıksal bir \textbf{Düğüm (Node)} hem de bu düğüm üzerinde çalışan birincil \textbf{Servis (Service)} olarak tanımlanır:

\begin{itemize}
    \item \textbf{Mantıksal Düğüm (Node ID):} \code{DEOS\_NODE\_BRAKE} (\hex{12})
    \item \textbf{Birincil Fonksiyonel Servis (Service ID):} \code{DEOS\_SERVICE\_BRAKE} (\hex{03})
    \item \textbf{Desteklenen Ortak Servisler:} \code{DEOS\_SERVICE\_SYSTEM} (\hex{00}), \code{DEOS\_SERVICE\_DIAGNOSTIC} (\hex{02})
\end{itemize}

\subsection{Gerçek Zamanlı Kontrol Komutları (Real-Time RX)}

\begin{center}
\begin{tabularx}{\textwidth}{>{\raggedright\arraybackslash}p{3.8cm} >{\raggedright\arraybackslash}p{3cm} >{\raggedright\arraybackslash}p{2.2cm} >{\raggedright\arraybackslash}p{2.2cm} X}
\toprule
\textbf{Değişken Adı} & \textbf{Sınıf \& Servis} & \textbf{Komut ID} & \textbf{Ölçek} & \textbf{Fonksiyonel Görev} \\
\midrule
\code{target\_brake\_demand} & COMMAND (0x1) BRAKE (0x03) & 0x01 & 0.01 \% & Operatör pedalından veya üst kontrolcüden ham fren talebi. \\
\code{target\_deceleration} & COMMAND (0x1) BRAKE (0x03) & 0x01 / 0x03 & 0.01 m/s² & Otonom planlayıcıdan gelen kinematik hedef yavaşlama ivmesi. \\
\code{brake\_control\_mode} & COMMAND (0x1) BRAKE (0x03) & 0x04 & Enum & Fren komutunun hangi modda icra edileceğini belirler (Yüzde, İvme, Basınç). \\
\code{target\_park\_brake} & COMMAND (0x1) BRAKE (0x03) & 0x05 & Enum & Elektronik Park Freni (EPB) ve rampada dur-kalk kilit desteği. \\
\code{brake\_prefill\_cmd} & COMMAND (0x1) BRAKE (0x03) & 0x06 & Boolean & Balata boşluğunu diske sıfırlayarak reaksiyon süresini kısaltır. \\
\code{estop\_brake\_override} & SAFETY (0x0) BRAKE (0x03) & 0x03 & Boolean & Freni tam sıkar; Kelly Pin 10'u opto ile toprağa çekip gazı keser. \\
\bottomrule
\end{tabularx}
\end{center}

\subsection{Fren Durum Telemetrisi (STATUS)}

\code{DEOS\_CMD\_BRAKE\_GET\_STATUS (0x02)} komutuna karşılık dönen veya periyodik yayınlanan veri yapısı:

\begin{center}
\begin{tabular}{c c c l}
\toprule
\textbf{Offset} & \textbf{Boyut} & \textbf{Tip} & \textbf{Alan} \\
\midrule
0 & 2 byte & uint16\_t & Current Position (0.01 mm) \\
2 & 2 byte & uint16\_t & Actuator Force (1 N) \\
4 & 2 byte & uint16\_t & Motor Current (0.1 A) \\
6 & 2 byte & int16\_t  & Chassis Long. Accel (0.01 m/s$^2$) \\
8 & 1 byte & uint8\_t  & Brake System Active (0: Idle, 1: Active) \\
9 & 1 byte & uint8\_t  & Park Brake State (0: Released, 1: Engaged) \\
\bottomrule
\end{tabular}
\end{center}

\subsection{Çapraz Sistem Dinamik Telemetrisi (Emniyet ve Kayma Algılama)}
\begin{itemize}
    \item \textbf{current\_vehicle\_speed} (0x10): Araç durmadan mekanik park kilidinin atmasını donanımsal olarak engeller.
    \item \textbf{traction\_wheel\_slip\_state} (0x10): Kilitlenme sezildiğinde fren aktüatörüne kaba ABS gevşetmesi uygular.
    \item \textbf{chassis\_longitudinal\_accel} (0x03): Fren kontrol döngüsünde ivme kapalı çevrim geri besleme girdisi sağlar.
\end{itemize}

\subsection{Fren Konfigürasyon Parametreleri}
\begin{center}
\begin{tabularx}{\textwidth}{>{\raggedright\arraybackslash}p{3.8cm} >{\raggedright\arraybackslash}p{3cm} >{\raggedright\arraybackslash}p{2.2cm} >{\raggedright\arraybackslash}p{2.2cm} X}
\toprule
\textbf{Parametre} & \textbf{ID} & \textbf{Tip} & \textbf{Ölçek} & \textbf{Açıklama} \\
\midrule
PID\_KP/KI/KD & \hex{0001-03} & uint16\_t & 0.01 & Fren pozisyon/basınç kapalı çevrim kazançları. \\
MIN\_POSITION & \hex{0010} & uint16\_t & 0.01 mm & Balatanın mekanik tam serbest kalma sıfır konumu. \\
MAX\_POSITION & \hex{0011} & uint16\_t & 0.01 mm & Balatanın azami basma strok sınırı. \\
CURRENT\_LIMIT & \hex{0012} & uint16\_t & 0.1 A & Aktüatör motorunun termal koruma akım tavanı. \\
SLEW\_RATE\_LIMIT & \hex{0013} & uint16\_t & 0.01 \%/ms & Mekanik dişli şoklarını önleyen maksimum basma hızı. \\
CMD\_TIMEOUT\_MS & \hex{0014} & uint16\_t & 1 ms & İletişim koptuğunda acil duruş frenine geçme gecikmesi. \\
KISS\_POINT\_OFFSET & \hex{0015} & uint16\_t & 0.01 mm/mV & Balatanın diske temas ettiği boşluksuz sıfır noktası. \\
MAX\_FORCE\_LIMIT & \hex{0016} & uint16\_t & 1N / 0.1bar & Mekanik aşırı yük koruma kuvvet tavanı. \\
\bottomrule
\end{tabularx}
\end{center}

\section{Kalibrasyon Prosedürleri}
Fren sistemi \code{DEOS\_SERVICE\_CALIBRATION (0x05)} servisini kullanır.
\begin{itemize}
    \item \textbf{Fren:} Balata temas noktası akım sıçraması okuması ile tespit edilir. Kaliper açıkken dara (Tare) alınır.
\end{itemize}

\end{document}
